\chapter{Analyzing a simulation}
\label{ch:analysis}

A simulation produces two kinds of result: a record of energies and
temperature, and a trajectory of snapshots showing where every atom was. The
analysis tools of the Molecular Dynamics Manager turn these into numbers you
can interpret: how far the protein moved, which parts are flexible, whether
it kept its shape, and much more. This chapter shows how to use them and how
to read the main results.

\begin{glance}
\glanceitem{What it does}{Checks the energies of each stage and measures the
  protein's motion and shape across the snapshots (23 kinds of
  measurement).}
\glanceitem{You need it when}{A simulation stage has finished (or to look at
  any Amber trajectory you have).}
\glanceitem{Before you start}{A finished stage from \key{5b}
  (Chapter~\ref{ch:runningmd}), or a trajectory (\file{.nc}) with its
  topology (\file{.prmtop}).}
\glanceitem{Where to find it}{\key{5b} \key{4} (Analyze completed
  simulations), \key{4} in the Execute Simulations menu, or the Terminal
  command \texttt{proprep --analysis}.}
\glanceitem{You get}{Results and text plots on screen, and on request a
  \file{.csv} and \file{.json} file of the numbers.}
\end{glance}

\section{Two ways in}

\begin{itemize}
  \item \textbf{From the menus:} \key{5b}, then \key{4} (Analyze completed
        simulations). After a run, the Execute Simulations menu offers the
        same tool as \key{4}.
  \item \textbf{Straight from the Terminal:} move into the project folder and
        type the command below. ProPrep skips the menus and opens the analysis
        directly.
\end{itemize}

\begin{terminal}
cd ~/proprep-tutorial/ubq
proprep --analysis --full-menu
\end{terminal}

Without \texttt{--full-menu}, ProPrep first asks \prompt{Select menu mode
[1/2]} and \prompt{Save this as your default preference? (n)}; answer \key{2}
and \key{n}. The command cannot be combined with \texttt{--pdbid} or
\texttt{--pdbfile}.

\section{Step by step: analyzing the ubiquitin run}\label{sec:analysis-steps}

\subsection{Choose the simulation}

ProPrep lists the finished stages in the project, newest first:

\begin{screen}
===== Simulation Analysis =====

Available simulations for analysis:
  1. 03_1UBQ_1UBQ_NVT_Production_MD (completed: 2026-11-21 12:58)
  2. 02_1UBQ_1UBQ_System_Heating (completed: 2026-11-21 12:58)
  3. 01_1UBQ_1UBQ_Energy_Minimization (completed: 2026-11-21 12:57)
Select simulation (1-3) or 'browse' (1): §ans«1»

Analyzing Simulation: 03_1UBQ_1UBQ_NVT_Production_MD

Available Analysis Files:
  • 1 mdout file(s) (energetics)
  • 1 trajectory file(s) (3.6 MB total)
  • 1 topology file(s)

Select Analysis Type:
  1. Energetics analysis (mdout)
  2. Trajectory analysis (structural)
  3. Combined analysis (recommended)
  4. View trajectory in the structure viewer
Select analysis type (3):
\end{screen}

\begin{steps}
  \item Type the number of the stage, or \typed{browse} to pick files
        anywhere (Section~\ref{sec:analysis-join}).
  \item Choose the type of analysis: \key{1} energies only, \key{2} the
        trajectory only, \key{3} both one after the other. A minimization has
        no trajectory, so ProPrep goes straight to the energies.
\end{steps}

Option \key{4} opens the trajectory in ProPrep's 3D viewer in your web browser
(Chapter~\ref{ch:compare}); it was not tried for this manual.

\subsection{Energetics}\label{sec:analysis-energy}

The \term{mdout} file is Amber's running report of energy, temperature and
pressure. ProPrep first guesses what kind of stage it was and asks you to
confirm:

\begin{screen}
Simulation Type Detection
System size: 15,729 atoms

Detected: NVT Equilibration/Production (confidence: medium)
Evidence:
  • ntp=0, ntb=1 (NVT ensemble)

Simulation Types:
  1. Minimization
  2. Heating (NVT)
  3. Heating (NPT)
  4. NVT Equilibration/Production (detected)
  5. NPT Equilibration/Production
Confirm simulation type (4): §ans«»
\end{screen}

The type decides which checks you get:

\begin{tabularx}{\linewidth}{@{}lL@{}}
\toprule
\textbf{Stage} & \textbf{What ProPrep reported for ubiquitin} \\
\midrule
Minimization & 300 steps; energy from $-26{,}515$ to $-60{,}537$\,kcal/mol;
  final RMS gradient 1.12; \prompt{Monotonic Decrease: Yes}. Plots of energy
  and gradient. \\
Heating & 2.00\,ps; 0.0\,K to 291.7\,K; average of the last 20\,\%
  282.7 ± 7.0\,K (target 300\,K). \\
Production & Temperature 294.50 ± 3.71\,K; total, kinetic and potential
  energy; plots of temperature and total energy over time and a temperature
  histogram. \\
\bottomrule
\end{tabularx}

\medskip
What to look for: in a minimization the energy should fall and level off; in
heating the temperature should reach the target; in production the
temperature should stay near the target and the energy should not drift
steadily up or down.

\begin{watchout}
In version 1.23.0 the heating stage was detected as \prompt{NVT
Equilibration/Production}. Check the detected type and type the right number
(\key{2} for heating at constant volume) before pressing \key{Enter}.
\end{watchout}

\subsection{Load the trajectory}

For \key{2} or \key{3}, a small file browser opens in the stage folder:

\begin{screen}
Select trajectory file(s) to analyze
...
  [ 1] .. (parent)
  [ 2] 1ubq_1ubq_nvt_production_md.nc  (3.6 MB, 11/21/2026 12:58)

Commands:
  N          - open dir N, or select file N
  N,M  N-M   - select multiple files (e.g. 1-3,5)
  ..         - go up one level
  find      - Recursive search for .nc files
  q          - cancel

Selection (q): §ans«2»
\end{screen}

Type the file's number (here \key{2}), or \typed{find} and then its number in
the search list. ProPrep loads it and summarizes:

\begin{screen}
        Trajectory Information
┌────────────┬───────────────────────┐
│ Frames     │ 20                    │
│ Atoms      │ 15,729                │
│ Residues   │ 4,926                 │
│ Time range │ 2.0 - 3.0 ps (0.0 ns) │
│ Box type   │ periodic              │
└────────────┴───────────────────────┘
\end{screen}

A \term{frame} is one snapshot. The residue count includes every water and
ion.

\subsection{Choose the analyses}

\begin{screen}
══ STRUCTURAL ANALYSIS ══
   1. RMSD - Root mean square deviation
   2. RMSF - Per-residue fluctuations
   3. Contacts - Contact maps & native contacts
   4. Salt bridges - Electrostatic interactions
   5. SASA - Solvent accessible surface area
══ SECONDARY STRUCTURE & GEOMETRY ══
   6. DSSP - Secondary structure (helix, sheet, coil)
   7. Ramachandran - Backbone dihedral angles (φ, ψ)
   8. B-factors - Pseudo B-factors from MD
══ PAIRWISE MEASUREMENTS ══
   9. Distance - Atom-atom distances
  10. Angle - Three-atom angles
  11. Dihedral - Four-atom torsion angles
  12. Vector - Orientation tracking
══ DYNAMICS & CORRELATION ══
  13. PCA - Principal component analysis
  14. Clustering - Conformational clustering
  15. Autocorrelation - Temporal correlations
  16. Pairwise RMSD - All-vs-all RMSD matrix
══ SOLVATION ANALYSIS ══
  17. Hydrogen bonds - H-bond network
  18. Water RDF - Radial distribution function
  19. Water shells - Hydration layer analysis
  20. Density maps - Spatial density distributions
══ GEOMETRIC PROPERTIES ══
  21. Radius of gyration - Structural compactness
  22. Contact frequency - Per-residue contact analysis
  23. Periodic image distance - Solute to its own copies across the box

   0. Exit to MD Manager menu

Select analyses (comma-separated, e.g., 1,2,6):
\end{screen}

Type one number, or several separated by commas, such as \typed{2,9,21,6}.
ProPrep runs them in that order (\prompt{Analysis 1/4}, \prompt{Analysis 2/4},
\ldots), asks a few questions for each, and then offers to save the numbers.

\begin{watchout}
In version 1.23.0, any answer that contains the digit \texttt{0} is treated
as \key{0} (Exit). Typing \typed{10}, \typed{20}, or a list such as
\typed{15,18,19,20,22}, prints \prompt{Returning to MD Manager menu\ldots} at
once. As a result, \textbf{Angle (10)} and \textbf{Density maps (20)} cannot
be run at all. Leave them out of your lists.
\end{watchout}

\subsection{Choosing the atoms: regions}

Most analyses first ask which atoms to use:

\begin{screen}
Common Selections:
  1. Backbone atoms (@C,CA,N,O&!:WAT)
  2. C-alpha atoms (@CA)
  3. All atoms (*)
  4. Protein residues 1-76 (no H, excludes solvent)
  5. Specific residues (custom)
  6. Custom AMBER mask (advanced)
Select region (2):
\end{screen}

The default, \key{2} (C-alpha atoms), is one atom per amino acid and suits
RMSD, RMSF and PCA. Choose \key{4} (the whole protein without hydrogens) for
size and shape measures such as radius of gyration, DSSP and surface area.
\key{5} and \key{6} accept residue ranges or an Amber mask
(Section~\ref{sec:runningmd-restraints} explains masks).

\subsection{Saving the numbers}\label{sec:analysis-export}

After each batch ProPrep asks:

\begin{screen}
Export Options:
  1. CSV format (for Excel, Origin, plotting)
  2. JSON format (for Python, R)
  3. Both formats
  4. Don't export
Export analysis results? (4): §ans«3»
✓ Exported to .../ubq/03_1UBQ_1UBQ_NVT_Production_MD_trajectory_analysis.csv
✓ Exported to .../ubq/03_1UBQ_1UBQ_NVT_Production_MD_trajectory_analysis.json
\end{screen}

The default is \key{4} (no file). The files go into the project folder and
are named after the stage. The \file{.csv} file has one row per frame and one
column per result (for example \texttt{rmsd\_Backbone\_atoms},
\texttt{rgyr\_\ldots}, \texttt{dssp\_helix\_pct}); open it in Excel or any
spreadsheet program to make your own graphs. The \file{.json} file holds the
same data for scripts.

\begin{tip}
Every export in the same session writes to the \emph{same} file names and
replaces the previous files. The new file contains everything analyzed so
far in the session, so the last export is the complete one. Copy or rename
the files if you want to keep a particular version.
\end{tip}

The plots appear only on screen, drawn with text characters. ProPrep does not
save picture files; make publication graphs from the \file{.csv} file.

\section{The analyses}\label{sec:analysis-list}

The table lists all 23. ``Tried'' shows what we found for the 1\,ps
ubiquitin production run (20 frames). Such a short run is only a test, so the
numbers show what the output looks like, not real biology.

\begin{longtable}{@{}>{\raggedright\arraybackslash}p{0.25\linewidth}>{\raggedright\arraybackslash}p{0.38\linewidth}>{\raggedright\arraybackslash}p{0.31\linewidth}@{}}
\toprule
\textbf{Analysis} & \textbf{Tells you} & \textbf{Tried on 1UBQ} \\
\midrule
\endhead
1 RMSD & How far the structure has moved from a reference frame &
  Backbone, frame 0: final 0.65\,Å \\
2 RMSF & How much each residue wobbles about its average position &
  C-alpha: mean 0.28\,Å; residue 76 highest (0.72\,Å) \\
3 Contacts & Which residue pairs touch; fraction of the starting
  (``native'') contacts kept & 162 native contacts, 158 persistent \\
4 Salt bridges & Pairs of opposite charges close together, and how often &
  8 salt bridges, all persistent \\
5 SASA & Surface area reachable by water & 2,940\,Ų on C-alpha atoms only, which is not meaningful; use region~4 \\
6 DSSP & Helix, sheet, turn and coil content over time &
  19.5\,\% helix, 33.9\,\% sheet, 19.0\,\% turn, 27.6\,\% coil \\
7 Ramachandran & Backbone angles φ and ψ (or χ1, ω) of each residue &
  23\,\% of residues in the α region, 31\,\% in β \\
8 B-factors & Flexibility expressed like crystallographic B-factors &
  Mean 2.22\,Ų \\
9 Distance & Distance between two atoms or groups over time &
  \texttt{:1@CA} to \texttt{:76@CA}: 38.44 ± 0.70\,Å \\
10 Angle & Angle between three atoms & Cannot be selected (see above) \\
11 Dihedral & Twist angle defined by four atoms &
  Phe4 χ1: mean $-51.7$° \\
12 Vector & Length and direction of a vector between two groups &
  8.43\,Å, polar angle 134.9° \\
13 PCA & The main directions of collective motion &
  PC1 54.7\,\%; first three 83.2\,\% \\
14 Clustering & Groups of similar structures &
  5 clusters (K-means) \\
15 Autocorrelation & How quickly a quantity ``forgets'' its earlier value &
  Ran on an RMSD series \\
16 Pairwise RMSD & RMSD between every pair of frames & 190 pairs, mean 0.40\,Å \\
17 Hydrogen bonds & Which hydrogen bonds exist and how often &
  36 ± 4 per frame; Asp52--Lys27 present 100\,\% \\
18 Water RDF & How water is arranged around a chosen group &
  First peak at 2.75\,Å, but see Section~\ref{sec:analysis-problems} \\
19 Water shells & Number of waters in layers around the protein &
  5 shells of 2\,Å \\
20 Density maps & Where solvent or ions spend time & Cannot be selected \\
21 Radius of gyration & Overall size and compactness &
  11.90 ± 0.03\,Å (protein, no H) \\
22 Contact frequency & How often each residue touches others & Prompts seen \\
23 Periodic image distance & Closest approach of the protein to its own copy
  & Same as monitor option 12 (Chapter~\ref{ch:runningmd}) \\
\bottomrule
\end{longtable}

Most analyzes ask a few more questions, each with a sensible default:

\begin{itemize}
  \item \textbf{RMSD:} region and reference frame (frame 0, the average
        structure, or a frame you name); then \prompt{Analyze additional
        regions?}.
  \item \textbf{RMSF:} region, the atoms used to line up the frames (default
        \texttt{@CA,C,N}) and the reference frame.
  \item \textbf{Distance:} two masks, such as \typed{:1@CA} and
        \typed{:76@CA}, and an optional label.
  \item \textbf{Contacts:} contact type, cutoff (4.5\,Å), how often a contact
        must be present to count (0.5 = half of the frames), and which frame
        defines the native contacts.
  \item \textbf{Salt bridges:} cutoff 4.0\,Å (or relaxed 5.0, strict 3.2).
  \item \textbf{Hydrogen bonds:} donor and acceptor masks (blank = all),
        distance 3.0\,Å and angle 135°.
  \item \textbf{PCA:} region, number of components, and whether to fit
        frames to the average first.
  \item \textbf{Clustering:} region, method (K-means or hierarchical), number
        of clusters, and the similarity measure.
  \item \textbf{Water shells:} protein or custom mask, shell width (2.0\,Å)
        and maximum distance (10\,Å).
\end{itemize}

\section{Reading RMSD, RMSF and radius of gyration}\label{sec:analysis-reading}

\subsection{RMSD: has the structure settled?}

The \term{RMSD} (root mean square deviation) compares each frame with a
reference frame, atom by atom, after lining the two up. It is an average
distance in Å. It starts at 0 (the reference compared with itself) and rises
as the protein moves away from its starting shape.

\begin{screen}
RMSD Analysis: Backbone atoms
Mask: @C,CA,N,O&!:WAT
Reference: Frame 0

     RMSD Statistics
┌───────┬───────────────┐
│ Mean  │ 0.50 ± 0.18 Å │
│ Range │ 0.00 - 0.65 Å │
│ Final │ 0.65 Å        │
└───────┴───────────────┘
\end{screen}

How to read it:
\begin{itemize}
  \item For a stable folded protein the backbone RMSD usually rises over the
        first part of the run and then levels off, often between 1 and 3\,Å.
        A plateau means the structure has settled.
  \item An RMSD that keeps climbing means the protein is still changing:
        the run may be too short, or the protein may be unfolding.
  \item A sudden jump points to a specific event, such as a loop or domain
        moving. Look at the trajectory around that time.
\end{itemize}

\subsection{RMSF: which parts are flexible?}

The \term{RMSF} (root mean square fluctuation) is measured per residue. It
says how far each residue wanders around its own average position during the
run.

\begin{screen}
     RMSF Statistics
┌────────┬───────────────┐
│ Mean   │ 0.28 ± 0.08 Å │
│ Range  │ 0.16 - 0.72 Å │
│ Median │ 0.27 Å        │
└────────┴───────────────┘

High Fluctuation Residues (top 10%, n=8):
  Residue 10: 0.35 Å
  Residue 15: 0.52 Å
...
  Residue 76: 0.72 Å
\end{screen}

Peaks in the RMSF plot are flexible regions: loops and chain ends. Low values
mark the rigid core, such as helices and sheets. In ubiquitin the C-terminal
residue 76 moved most, as expected for a free chain end. Compare RMSF with
the crystallographic B-factors of the PDB file, or use analysis~\key{8}, which
converts RMSF to the same scale ($B = 8\pi^2/3 \cdot \mathrm{RMSF}^2$).

\subsection{Radius of gyration: is it still compact?}

The \term{radius of gyration} (Rg) is the average distance of the atoms from
the protein's center. A compact, folded protein has a small, steady Rg;
unfolding makes it grow. Ubiquitin's Rg stayed at 11.90 ± 0.03\,Å. A steady
Rg together with a levelled-off RMSD is good evidence that the fold is
stable.

\subsection{The other common measures}

\begin{itemize}
  \item \textbf{DSSP} assigns each residue to helix, sheet, turn or coil in
        every frame. Steady percentages mean the secondary structure is kept.
  \item \textbf{Hydrogen bonds} and \textbf{salt bridges} report how often
        each interaction is present (in \% of frames). Interactions present
        close to 100\,\% are stable features of the structure.
  \item \textbf{PCA} (principal component analysis) finds the few directions
        that describe most of the motion. A large first component (here
        54.7\,\%) means one motion dominates.
\end{itemize}

\section{Joining trajectories with \texttt{proprep --analysis}}\label{sec:analysis-join}

To analyze several stages as one continuous run (for example heating
followed by production), type \typed{browse} at \prompt{Select simulation}.
We did this after starting with \texttt{proprep --analysis}; the
\prompt{Select simulation} screen is the same one you reach with \key{5b}
\key{4}. The file browser that opens lets you pick files from different
folders:

\begin{tabularx}{\linewidth}{@{}lL@{}}
\toprule
\textbf{Command} & \textbf{What it does} \\
\midrule
\typed{select 3} or \typed{select 3,5} & Pick files and continue at once \\
\typed{add 3} & Put a file aside and keep browsing \\
\typed{done} & Continue with the files put aside \\
\typed{clear} & Forget the files put aside \\
\typed{cd 2}, \typed{cd ..}, \typed{cd} \textit{path} & Change folder \\
\typed{find} & Search all subfolders for analysis files \\
\typed{q} & Cancel \\
\bottomrule
\end{tabularx}

\medskip
In our test, \typed{find} listed 16 files (energy reports, trajectories and
topologies); we typed \typed{5,4,14} for the heating trajectory, the
production trajectory and the topology in \file{shared/}:

\begin{screen}
Selected 3 file(s)

Using topology: 1UBQ.prmtop

Loading 2 Trajectory Segments
  1. 1ubq_1ubq_system_heating.nc (3.6 MB)
  2. 1ubq_1ubq_nvt_production_md.nc (3.6 MB)

Concatenating segments...
✓ Loaded 40 total frames
\end{screen}

List the trajectories in time order. The radius of gyration over all 40
frames was 11.84 ± 0.06\,Å.

\begin{tip}
The search list may include empty topology files left by failed tLEaP runs
(shown as \prompt{0.0 KB}). Pick a topology of normal size; the copy in the
batch's \file{shared/} folder is a safe choice.
\end{tip}

The simpler browser reached through \key{5b} \key{4} \key{2} has a
\typed{find} that searches all stage folders, but there it accepts only
\emph{one} file; \typed{2,1} was rejected. To join files, use \typed{browse}
as described above.

Typing \key{0} in the analysis menu started from \texttt{proprep --analysis}
takes you to the ProPrep main menu, not back to the list of simulations.

\section{Other choices}

\subsection{Plugins}

ProPrep can be extended with \term{plugins}: separately installed add-on
packages that add their own items to ProPrep's menus. ProPrep finds them
automatically in the environment it runs in. No plugin was installed for this
manual, and the 23 trajectory analyzes are all built in.

\subsection{PB Titrate}

The Topology Generator's option \key{5} (Refine protonation states via PB)
calculates pKa values from the structure rather than analyzing a trajectory.
It is described briefly in Chapter~\ref{ch:topology}; its calculations were
not run for this manual.

\section{If something goes wrong}\label{sec:analysis-problems}

\begin{tabularx}{\linewidth}{@{}lL@{}}
\toprule
\textbf{What you see} & \textbf{Why, and what to do} \\
\midrule
\prompt{Returning to MD Manager menu\ldots} right after you list analyses &
  Your answer contained a 0 (for example 10 or 20). Choose again without
  them. Angle and Density maps cannot be run in version 1.23.0. \\
\prompt{Failed to setup action} and a long error report, then the MD Manager
  menu &
  A mask did not match any atoms; in our test a dihedral used atom
  \texttt{CB} of glycine 10, which has none. The loaded trajectory is lost;
  start the analysis again with valid atoms (for example
  \typed{:4@N :4@CA :4@CB :4@CG}). \\
Water RDF with an empty \prompt{Solute:} &
  The solute mask was left blank and ProPrep accepted it. Type a mask such as
  \typed{:1-76}. \\
SASA gives a large, meaningless number &
  The default region (C-alpha atoms) is a poor choice for surface area. Use
  region \key{4}. \\
Heating reported as production &
  Correct the detected type when asked (Section~\ref{sec:analysis-energy}). \\
The \file{.csv} lost an earlier result &
  Later exports overwrite it; the last one holds everything from the
  session. \\
\prompt{Select file} refuses \typed{2,1} &
  That browser takes one file. Type \typed{browse} at \prompt{Select
  simulation} instead (Section~\ref{sec:analysis-join}). \\
Lines starting with a date and \texttt{INFO} appear &
  Internal log messages; they can be ignored. \\
\bottomrule
\end{tabularx}

\medskip
See also Chapters~\ref{ch:troubleshooting} and~\ref{ch:knownissues}.
